\section{Introduction}

FragPicker~\cite{park2021fragpicker} avoids the cost of defragmentation on modern
storage by moving only the pieces of a file whose requests straddle extent boundaries.
Its artifact first traces an application with BCC and then runs an \emph{analysis
phase}, a chain of Python and shell scripts that turns the trace into per-file lists of
hot regions, before a \emph{migration phase} rewrites them.

\forkowner~\cite{forkbranch} parallelised the analysis phase: one batch inode-to-path
map, extent lookups through the FIEMAP ioctl, and a pool of 16 threads. It reported a
12.74$\times$ speedup when a 5\,ms \texttt{time.sleep} stood in for a slow disk before
every extent lookup; no slow device was measured. We ask whether the result holds and
find that it does not, for reasons that apply to any tool that does per-file metadata
work in Python:
\begin{itemize}
\item \textbf{A single thread beats every thread count.} A rewrite with byte-identical
  output runs the original's \texttt{processing.py} 4--191$\times$ faster on five real
  traces and is 19--26$\times$ faster than the best parallel setting (\S\ref{ws:time}).
\item \textbf{Sleeping is not I/O.} The user-space delay reproduces the published
  speedup (15.5$\times$); the same delay at the block layer leaves almost nothing to
  overlap, and threads make the lookups slower (\S\ref{ws:models}).
\item \textbf{The interpreter decides.} When lookups do wait, CPython 3.12 holds its
  global interpreter lock through the fork's 1{,}824-byte FIEMAP request, so 16 threads
  gain nothing; with 984 bytes, or under CPython 3.14, they gain 13--15$\times$. The
  threshold is a step at 1{,}024 bytes, and the same lock caps how fast threads start
  the parallel path's \texttt{sort} processes (3.5$\times$), yet removing it does not
  make the parallel path competitive (\S\ref{ws:gil}).
\item \textbf{Threads belong in the tree walk} (exploratory). In a cold, scattered tree
  the device time is in the walk that precedes the threads; a walk split across threads
  is 3.9$\times$ faster under emulated delay and 6.1$\times$ on a real network disk
  (\S\ref{ws:walk}).
\end{itemize}
The parallel mode also computes extents that nothing reads, lists no files at all on
file systems without FIEMAP, and deadlocks above 10{,}000 traced files (\S\ref{ws:correct}).
All scripts, traces and raw data are released~\cite{studybranch}.

\section{The analysis phase and the fork}
\label{ws:bg}

Upstream's \texttt{processing.py} has two loops. Phase~1 reads the parsed trace line by
line and appends each request to \texttt{<inode>.txt}, opening and closing the file for
every line. Phase~2 visits every traced inode, runs \texttt{find /mnt -inum} to turn it
into a path, sorts the request file and renames it. \texttt{merge.py} and
\texttt{hotness.sh} then merge adjacent requests and keep the hottest. Counted with
\texttt{strace}, this costs fourteen new processes per traced file (3 in
\texttt{processing.py}, 1 in \texttt{merge.py}, 10 in \texttt{hotness.sh}). The
analysis never reads extents: only the migration phase does, with \texttt{filefrag}.

The fork's \texttt{--parallel} path builds the inode map with one \texttt{find}, gives
each traced file to a thread that runs FIEMAP on it, writes the extents to a temporary
directory that nothing reads, and then sorts the request files.

\section{Method}
\label{ws:method}

\paragraph{Ablation ladder.} (a) upstream unmodified; (b) the fork's sequential path;
(c) our single-threaded rewrite (one batch inode map, Phase~1 in memory, no
subprocesses), whose output is byte-identical to (a); (d) \texttt{--parallel} with 1, 4,
8 or 16 threads; (e) (d) with the extent lookup replaced by a stub; and (m) (d) with the
fork's resource-monitor thread replaced by a no-op. Every run hashes its output: (c)
matched (a) in every run, and the fork's rungs matched each other everywhere and matched
(a) except on traces with requests not aligned to blocks (\S\ref{ws:correct}).

\paragraph{Workloads.} Five real traces captured with the artifact's own tracer
(upstream's two read benchmarks, \texttt{tar}, \texttt{fio}, \texttt{sqlite3}), and
synthetic traces of $F$ files chosen from trees of $T$ files.

\paragraph{Slow devices.} The fork's \texttt{sleep}; dm-delay~\cite{kerneldmdelay}
below the file system, which delays only reads that reach the device and lets any
number wait at once; and an NBD device served by nbdkit~\cite{nbdkit} with one thread
(one request at a time) or 32. A configured 5\,ms gives 8.0\,ms per read under dm-delay
(its timer counts in 4\,ms jiffies) and 5.6\,ms under NBD, measured with \texttt{fio}.
Caches are dropped and the file system remounted before every cold run.

\paragraph{Platform and statistics.} One laptop (Ryzen~5 6600H, 6 cores, 15\,GB, NVMe
SSD) running Linux 6.18 in WSL2 with CPython 3.12.3; test file systems (ext4, XFS,
Btrfs, F2FS) on loop devices with direct I/O. Each configuration runs 5--10 times in
shuffled order; we report medians with 95\% bootstrap intervals~\cite{hoefler2015scientific}.
We wrote our predictions down before the first run; the experiments of
\S\ref{ws:walk} and the interpreter comparison of \S\ref{ws:gil} were designed after
seeing the data and are exploratory. One window of external load on the laptop was
detected after the fact and its workload repeated.

\section{Results}

\subsection{Does the parallel path compute the same thing?}
\label{ws:correct}
Only as a side effect: a file whose FIEMAP fails is dropped from \texttt{filelist.txt},
so on tmpfs and on WSL's 9p file system the parallel path lists no files where upstream
lists all of them. Its task queue is filled before the workers start, so the 10{,}001st
traced file blocks forever, and it re-sorts unrelated \texttt{*.txt} files in the working
directory. Its \texttt{analyze.sh} does not parse. Both of the fork's paths also round the
end of an unaligned request differently from upstream (neither formula is exact), so on
the buffered \texttt{tar} and \texttt{sqlite3} traces their hot regions differ from
upstream's.

The trace that every rung analyses is itself lossy, which affects upstream and the fork
alike. BCC's \texttt{trace} flushes its output only after the header, and upstream stops
it with SIGTERM, so the last block-buffered chunk of events is lost: 59 of the 8{,}192
reads of upstream's own \texttt{read\_seq} benchmark, and every event of a workload with a
single read. Unbuffered output recovers the tail but makes BCC drop samples at high
rates (3{,}739 during \texttt{tar}, about 13\%). File positions are printed with
\texttt{\%u}, so a read at 5\,GiB is recorded at 1\,GiB. Trace lines carry an inode number
but no device, so the 2 to 89 inodes per trace that live on the root file system
(shared libraries, locale files) are attributed to any file under \texttt{/mnt} that
happens to have the same number.

\subsection{Where the time goes}
\label{ws:time}

\begin{figure}[t]
\centering
\figfile{width=\columnwidth}{../figs/e2_vsF_ext4_T10000_warm_None.pdf}
\caption{Analysis time against the number of traced files $F$ in a tree of 10{,}000 files
(warm cache; medians, 95\% intervals).}
\label{ws:fig:ladder}
\end{figure}

Phase~1 makes about 41 system calls per trace line, 30 of them one-byte reads because
upstream opens the trace unbuffered; it costs 68--79\,\textmu s per line on all five file
systems. With 20 trace lines per file, \texttt{processing.py} costs 7.2\,ms per traced
file plus 0.80\,\textmu s per entry of the tree, a fit within 2\% at trees of 1{,}000,
10{,}000 and 100{,}000 files: every traced file pays for a walk of the whole tree.
Table~\ref{ws:tab:real} gives the five real traces. Where a workload touches many files
the per-file steps dominate: the whole pipeline takes 165\,s for \texttt{tar}'s 5{,}100
files, of which Phase~1 is 1.7\,s, against 1.28\,s for the rewrite (129$\times$). Where it
touches one file, Phase~1 dominates, and the parallel path is slower than upstream.

\begin{table*}[t]
\centering
\caption{Real traces ($L$ lines, $F$ traced files), warm cache: seconds, median of seven
runs, for upstream \texttt{processing.py} (a), the best \texttt{--parallel} setting (d),
the rewrite (c), and the whole pipeline (\texttt{processing.py}, \texttt{merge.py},
\texttt{hotness.sh}) against the rewrite doing all three steps.}
\label{ws:tab:real}
\small
\inputtable{../results/tab_real.tex}
\end{table*}

The parallel path never costs less than 1.27\,s (Figure~\ref{ws:fig:ladder}): its
resource-monitor thread sleeps for a second and is joined only after the sleep ends
(rung m: 0.22\,s). It is faster than upstream only because one walk replaces the
per-file \texttt{find}, and it still starts two \texttt{sort} processes per file. With 1{,}000
traced files the rewrite takes 0.19\,s, 80$\times$ faster than upstream and 25$\times$
faster than the best parallel setting (4.68\,s, 8 threads); stubbing out the extent
lookup saves the parallel path only 24\%. No thread count we tried beats the single
thread, so the parallel engine's COST~\cite{mcsherry2015cost} is unbounded.

\subsection{Sleeping is not I/O}
\label{ws:models}

\begin{figure*}[t]
\centering
\figfile{width=0.82\textwidth}{../figs/e5_models_d5.pdf}
\caption{The fork's lookup engine for 1{,}000 traced files under three models of a 5\,ms
device (``seq'': the plain sequential loop). Left: files created together; right: files
sampled from a tree of 100{,}000. Labels: data actually read from the device, cold.}
\label{ws:fig:models}
\end{figure*}

With \texttt{time.sleep} before every lookup, 16 threads are 15.5$\times$ faster than
the sequential loop (Figure~\ref{ws:fig:models}), close to the published 12.74$\times$.
With the delay at the block layer, the 1{,}000 files created together are looked up in
0.164\,s from a cold cache, because the whole job reads 192\,KiB: ext4 reads inode-table
blocks ahead in groups of 32, and one read brings in the inodes of 512 files. Threads
only add overhead (0.47$\times$), and on the one-request-at-a-time device as well
(0.43$\times$). When the 1{,}000 files are scattered over 100{,}000, a cold run reads
28.5\,MiB and takes 9.8\,s, and 16 threads take as long. Splitting the run shows why:
the walk that builds the inode map reads 29\,MiB, and the lookups after it read nothing.
In FragPicker's own workflow the analysis runs right after the traced application, so
these metadata are usually cached anyway.

\subsection{The interpreter lock}
\label{ws:gil}

\begin{table}[t]
\caption{Speedup of 16 threads over one for FIEMAP calls that each wait for one device
read (200 files with 16 extents, cold), by request size and interpreter; 3.14t has no
GIL. Medians of five passes.}
\label{ws:tab:gil}
\small
\setlength{\tabcolsep}{4pt}
\inputtable{../results/tab_gilpy_ws.tex}
\end{table}

Python threads overlap only outside the interpreter. In CPython 3.12 and 3.13,
\texttt{fcntl.ioctl} releases the GIL only when its argument fits a 1{,}024-byte stack
buffer; the fork's FIEMAP request holds 32 extents, 1{,}824 bytes, so the call holds the
lock while the kernel reads the extent tree. Table~\ref{ws:tab:gil} isolates this with
files whose extent block must be read: on 3.12 the fork's request gains nothing on any
device, the same request with room for 17 extents (984 bytes) gains 14.6--14.9$\times$ on
devices that serve requests concurrently, and the unchanged 1{,}824-byte request gains
14.6--14.9$\times$ under CPython 3.14, which releases the lock on this path (the 3.14
runs used a small stand-in for the two \texttt{psutil} calls the fork makes at import). Timing each
call separately confirms where the wait is: on ext4 and XFS the one read per fragmented
file happens inside FIEMAP; on F2FS it happens in \texttt{stat}, which releases the lock,
and there 16 threads speed that call up 3.5--4.7$\times$. The \texttt{sleep} of the
original evaluation released the lock and ran before the ioctl, so it hid this as well.

\begin{figure}[t]
\centering
\figfile{width=0.92\columnwidth}{../figs/e13_boundary.pdf}
\caption{16 threads over one for FIEMAP calls that each wait for one device read
(dm-delay), against the request size. Medians of three runs.}
\label{ws:fig:step}
\end{figure}

Sweeping the request from 1 to 64 extents turns the comparison into a step
(Figure~\ref{ws:fig:step}): on 3.12, 16 threads gain 14.5--14.9$\times$ up to 17 extents
(984 bytes) and 0.99--1.00$\times$ from 18 extents (1{,}040 bytes) on, and a sampler in a
separate process sees the device given one read at a time. CPython's rewrite of the
\texttt{fcntl} module removed the threshold in 3.14~\cite{cpython132742}. The same lock
caps the parallel path's main cost: from threads, \texttt{sort} processes start at most
3.5$\times$ faster than from one thread on 3.12 and 3.14, while worker processes, or
threads on the free-threaded 3.14t, reach 9--10$\times$. Removing the lock still does not
rescue the parallel path: on 3.14t it takes 3.43\,s against the rewrite's 0.27\,s.

\subsection{Where threads could pay}
\label{ws:walk}

\begin{figure}[t]
\centering
\figfile{width=0.92\columnwidth}{../figs/walk_threads.pdf}
\caption{A cold walk of 100{,}000 files split across threads, against one thread
(medians). The real network disk was measured at 1 and 16 threads only.}
\label{ws:fig:walk}
\end{figure}

Since a cold, scattered run spends its device time in the walk, we split an
\texttt{os.scandir} walk of the 100{,}000-file tree across threads by top-level
directory; \texttt{opendir} and \texttt{readdir} release the lock, and the inode map
equals \texttt{find}'s in every run (Figure~\ref{ws:fig:walk}). Under dm-delay one thread
takes 9.7\,s and four take 2.5\,s (3.9$\times$); sixteen take 4.4\,s, as passing the lock
between them starts to cost. On the free-threaded 3.14t, sixteen take 0.70\,s (14$\times$). With 32 requests in flight on NBD the walk drops from 7.1 to
1.8\,s; on the one-at-a-time device nothing gains more than 1.06$\times$; on the fast loop
device threads are slower. We repeated a subset on a virtual machine whose disk is
network block storage (CPython 3.14.4, Linux 7.0). Random reads there took 35--52\,ms at
the median, with about ten times more throughput at 16 in flight. There 16 threads walk
the tree 6.1$\times$ faster than one and 9.5$\times$ faster than \texttt{find} (1.17 against
11.1\,s), while for 1{,}000 traced files in a cold tree of 10{,}000 upstream still takes
11.6\,s, the best parallel setting 2.42\,s, and the rewrite 0.29\,s.

The kernel already overlaps much of this work. Under dm-delay, ext4's inode-table
readahead (32 blocks by default) makes the walk read 3.6 times more data than with
readahead off (7{,}287 reads against 2{,}023) and still finish 1.7$\times$ faster (9.8
against 16.2\,s), because the readahead reads are issued together; for files created
together, 16 threads never gained more than 1.08$\times$ at any setting from 0 to 256.

\section{Discussion}

\paragraph{Lessons.} The 12.74$\times$ came from the latency model, not from the code.
Three habits would have caught it: compare against the best sequential
implementation~\cite{mcsherry2015cost}, not the parallel engine on one thread; inject
latency below the cache, where only real device reads are slowed, rather than before
every call; and count device reads per unit of work, which tells in advance how much
there is to overlap. For a Python tool, a fourth: check which calls release the
interpreter lock, and on which interpreter versions.

\paragraph{What the analysis phase needs.} None of the fixes needs parallelism: build
Phase~1 in memory, resolve inodes with one walk (or, better, have the tracer record
device numbers and file handles, which also removes the foreign-inode problem), and look
extents up once, in migration, which reads them anyway. If the metadata are cold and
the device is deep, split the walk across a few threads.

\paragraph{Threats.} Our measurements are virtualised and, except for one virtual
machine on network block storage, on one NVMe SSD; the emulated devices bound device
concurrency from both sides but model neither seek distance nor reordering. The network
disk served other programs during the runs, cached recently written data, and was
measured with three repetitions. All rungs are Python, as the artifact is.

\section{Related work}

FragPicker~\cite{park2021fragpicker,park2023tocs} motivates selective defragmentation on
flash and Optane; its artifact evaluation~\cite{sosp21ae} reproduced the synthetic
workloads, and ScaleDefrag~\cite{lee2026scaledefrag} parallelises e4defrag's
fragmentation check and data relocation; neither measures the cost of the analysis
phase that is our subject. Benchmarking studies have shown how strongly storage results
depend on cache state and setup~\cite{traeger2008nine,tarasov2011benchmarking} and how
innocent-looking setup biases measurements~\cite{mytkowicz2009wrong}, and Hoefler and
Belli~\cite{hoefler2015scientific} catalogue reporting practice for parallel speedups,
which we follow. Aging studies~\cite{smith1997aging,conway2017aging,kadekodi2018geriatrix}
show how real file systems scatter metadata; our scattered layout only approximates
aging. David~\cite{agrawal2011david} emulates large storage systems by modelling their
runtime; dm-delay~\cite{kerneldmdelay} and nbdkit~\cite{nbdkit} inject latency below the
file system, and we use them as two bounds on device concurrency.

\section{Conclusion}

A thread pool that reproduced a 12.74$\times$ speedup against a user-space sleep loses
to one thread once delays sit where devices impose them. The waits that remain are in
the tree walk, where a few threads do help, and under CPython 3.12 the fork's FIEMAP
request kept even those lookups that wait from overlapping. Scripts, traces, raw data and
the full report are in the artifact~\cite{studybranch}.
